\section{引言}

跟踪器丢了一个人、又找回来的时候，必须决定这个回来的人是不是原来那个。ByteTrack
\cite{zhang2022bytetrack} 会把丢失的轨迹保留三十帧，之后任何越过关联阈值的检测都会被接回去。接对了，身份穿过一次遮挡继续存在；接错了，没有任何提示：错的人继承了那个编号，此后的每一帧都挂在这个编号下，基于它统计的人数、停留时间和轨迹从那一刻起全部是错的。重连错误是静默的，这就是它比它在关联事件中的占比更值得处理的原因。

现有工作走两条路。一条改进关联代价本身，靠运动模型、匹配级联和以观测为中心的更新
\cite{zhang2022bytetrack,aharon2022botsort,cao2023ocsort}；另一条引入外观，或在跟踪器内部加一个重识别分支，或在跟踪结束后做链接。本文问一个两者都没有直接回答、但更窄的问题：\textbf{面对跟踪器已经提出的重连提案，哪些应当被接受？}

我们把重连当作一次决策，在它前面放一个标量门控（gate）。门控读取检测置信度、预测框残差、轨迹时长、观测到的间断长度，以及该提案与同一行、同一列中最优竞争者在代价上的裕量。这些量进入一个带正则的逻辑回归评分，乘积项让运动一致性和分配歧义的权重随间断长度改变。被拒绝的检测以未匹配状态交回跟踪器，因此这个决定会改变后续帧看到的状态。门控不训练外观网络，也不保存特征库，只依据几何信息作判断。

本文给出四条主张。\textbf{第一，效果在接受量对齐后依然存在。} 只数错误会奖励"少接"，所以我们重新选阈值，让原生关联代价、加性评分和交互评分在每一来源上接受同样多的动作。在每来源最多接受十二次的设置下，ByteTrack 的 381 次可评估重连中有 98 次失败，原生代价是 390 次中 158 次失败；核验正确的重连从 232 次升到 283 次。\textbf{第二，对照对象是已发表的评分，不只是手工阈值。} 把 BoostTrack 的关联相似度移植进同一宿主并逐元素核对之后，交互评分在全部五个工作点上都更少出错，差距为 9.72 到 13.50 个百分点。\textbf{第三，增益来自乘积项，而不是新增的竞争摘要。} 只保留三个竞争摘要、去掉乘积项的评分，在五个工作点上没有任何一个对照区间排除零；完整评分则在效果成立的那个工作点上把失败比例降低 3.29 个百分点。\textbf{第四，代价是明确的。} 相对于不限额的 ByteTrack，门控减少了 1,397 次可评估失败，也少恢复 136 个自然缺口，HOTA 下降 1.65 分。我们把它作为决策的价格报告，而不是设法把它解释掉。

有一条边界贯穿全文：门控不使用外观信号，因此结论适用于只有检测器的流水线；文中每一个失败比例都是对所指划分的观测值，不是对未见来源的界。
